\documentclass[10pt,a4paper]{amsart}
\usepackage[margin=19mm]{geometry}
\usepackage{amsmath,amssymb,mathtools}
\usepackage[scr=rsfs,cal=euler]{mathalfa}
\usepackage{xfrac}
\usepackage[unicode,hidelinks,bookmarks=false]{hyperref}
\newcommand{\ie}{i.e.\ }
\newcommand{\cf}{cf.\ }
\newcommand{\resp}{respectively\ }
\newcommand{\Subsection}{\subsection}
\providecommand{\nobreakdash}{\nobreak\hbox{-}}
\setlength{\emergencystretch}{2em}
\multlinegap0pt
\usepackage{amssymb}
\newtheorem{thm}{Theorem}
\newtheorem{prop}{Proposition}
\title{Quaternionic M\"{o}bius invariant Laplacian and quaternionic M\"{o}bius harmonic functions on the unit ball}
\author{Ruiwen Wang}
\date{Source: arXiv:2608.02153v2 (CC BY 4.0)}
\begin{document}
\maketitle

\noindent\emph{Source and licence.} Quaternionic M\"{o}bius invariant Laplacian and quaternionic M\"{o}bius harmonic functions on the unit ball. Version arXiv:2608.02153v2 (CC BY 4.0), distributed by arXiv under the Creative Commons Attribution 4.0 International licence, http://creativecommons.org/licenses/by/4.0/, as recorded in the arXiv licence metadata for this exact version and shown on the abstract page https://arxiv.org/abs/2608.02153v2. Full licence notice: \texttt{licenses/arxiv-2608.02153-CC-BY-4.0.txt}.\par\smallskip

\noindent\textit{Editorial scope.} Counted unit: the article's thm posson (source0.tex lines 118--137). The article's complete original proof of it is delivered with it (source0.tex lines 1069--1144, one \texttt{proof} environment of 76 lines, which the article places away from the statement). No mathematics is written, repaired or re-derived: the statement and the proof are the article's own lines, in the article's own notation, and no retained line is substituted or re-flowed.  Retained uncounted as the source-local context the proof needs: lines 407--409; lines 873--875; lines 969--975; lines 1051--1055.  Nothing was omitted from any retained span, so the package carries no omission record.  No retained line contains a citation, so the wrapper carries no bibliography and no bibliography entry.  No retained line contains a figure environment, an image-inclusion command or drawing code, so no image is delivered and the package declares no asset dependency.  Editorial formatting only, in addition: an AMS \texttt{amsart} preamble that restates the article's own declarations and macros used by the retained lines, namely the environments \texttt{thm}, \texttt{prop} on separate counters, together with the standard packages those lines need.  The retained environments print in this document as Theorem 1, Proposition 1 in order of appearance, whereas the article prints the same items with the numbering of its own full document.  The complete native source of the article as distributed in the arXiv e-print for this exact version is bundled unchanged as \texttt{sources/206674/source0.tex}.
\begin{equation}\begin{aligned}\label{z jab}
z_A^{1'}=\sum_B-J_A^B\overline{z_B^{0'}},\qquad\sum_{A,B}z_A^{A'}J_A^Bz_B^{B'}=\varepsilon^{A'B'}|\mathbf{q}|^2,
\end{aligned}\end{equation}

\begin{equation}\begin{aligned}\label{g eq}
r^2(1-r^2)g''(r^2)+ \bigl(2n+a+b-(a+b)r^2\bigr)g'(r^2)- a(b-1)g(r^2)=0,\\
\end{aligned}\end{equation}

\begin{prop}\label{Green}
Let $u,v\in C^2(B^{4n})$ and $r<1$. Then,
\begin{equation}\begin{aligned}\label{eq green}
\int_{B_{r}^{4n}}(u\triangle v-v\triangle u) d\widetilde{V}=2\int_{S^{4n-1}}(uR^{0'}v-vR^{1'}u)(r\zeta)\frac{r^{4n-2}d\sigma(\zeta)}{(1-r^2)^{2n}}.
\end{aligned}\end{equation}

\end{prop}

\begin{thm}\label{max}
Suppose $\Omega$ is an open subset of $B^{4n}$ and
that $f \in C^2(\Omega)$ is $\mathcal{QM}$-subharmonic on $\Omega$ and continuous on $\overline{\Omega}$ . If $f \leq 0$ on $\partial \Omega$,
then $f \leq 0$ in $\Omega$.
\end{thm}

\begin{thm}\label{posson}
For $\varphi\in C(S^{4n-1})$, the Dirichlet problem
\begin{equation}\begin{aligned}\label{Dirichlet p}
\left\{
\begin{aligned}
 \triangle u&=0 {\rm\quad on}\quad B^{4n},\\
u&=\varphi {\rm \quad on}\quad S^{4n-1},
\end{aligned}
\right.
\end{aligned}\end{equation}
has a unique solution, which is given by the following $\mathcal{QM}$-Poisson integral
\begin{equation}\begin{aligned}\label{def poisson int}
P[\varphi](\mathbf{a}):=\int_{S^{4n-1}} P(\mathbf{a},\zeta)\varphi(\zeta)d\sigma(\zeta),
\end{aligned}\end{equation}
 where $\sigma$ is the ${\rm Sp}(n)$Sp$(1)$-invariant measure on $S^{4n-1}$ satisfying $\sigma(S^{4n-1})=1$, and
\begin{equation}\begin{aligned}\label{def poisson k}
P(\mathbf{a},\zeta):=\left(\frac{1-|\mathbf{a}|^2}{|1-\mathbf{a}^*\zeta|^2}\right)^{2n+1}
\end{aligned}\end{equation}
is the $\mathcal{QM}$-Poisson kernel for $\triangle$ on $B^{4n}\times S^{4n-1}$.
\end{thm}

\begin{proof}[Proof of Theorem \ref{posson}]
For $g\in C^2([0,1])$, the equation $\triangle g(|\mathbf{q}|^2)=0$ is equivalent to
\begin{equation}\begin{aligned}
0=r^2(1-r^2)g''(r^2)+2ng'(r^2),
\end{aligned}\end{equation}
by \eqref{g eq} with $a=b=0$ for radial function $g$. It has a solution
\begin{equation}\begin{aligned}\label{g}
g(t)=\int_{t}^{1}\left(\frac{1-s}{s}\right)^{2n}ds.
\end{aligned}\end{equation}
For fixed $\mathbf{a}\in B^{4n}_r$, let $\Omega=B_r^{4n} \setminus \varphi_{\mathbf{a}}(B^{4n}(\mathbf{0},\epsilon))$ with $\epsilon$ sufficiently small, and let
\begin{equation}\begin{aligned}\label{Ga}
G_{\mathbf{a}}(\mathbf{q})=g(|\varphi_{\mathbf{a}}(\mathbf{q})|^2),
\end{aligned}\end{equation}
which is smooth Green function on $B^{4n} \setminus \{\mathbf{a}\}$ with singularity at $\mathbf{0}$. We also denote $G(\mathbf{q}):=g(|\mathbf{q}|^2)$.
To apply Green formula, we choose an extension $\widehat{v}\in C^2(B_r^{4n})$ such that $\widehat{v}=G_{\mathbf{a}}$ on $\overline{\Omega}$.
Thus
\begin{equation}\begin{aligned}\label{total}
0=\int_{\Omega}(u\triangle \widehat{v}-\widehat{v}\triangle u) d\widetilde{V}=\int_{B^{4n}_r}(u\triangle \widehat{v}-\widehat{v}\triangle u) d\widetilde{V}-\int_{\varphi_{\mathbf{a}}(B^{4n}(\mathbf{a},\epsilon))}(u\triangle \widehat{v}-\widehat{v}\triangle u) d\widetilde{V}.
\end{aligned}\end{equation}
Applying the Green formula in Proposition \ref{Green} to $u\circ \varphi_{\mathbf{a}}$ and $\widehat{v}\circ \varphi_{\mathbf{a}}$ on $B^{4n}(\mathbf{0},\epsilon)$, we get
\begin{equation}\begin{aligned}\label{part ii}
\int_{\varphi_{\mathbf{a}}(B^{4n}(\mathbf{0},\epsilon))}(u\triangle\widehat{v}-&\widehat{v}\triangle u) d\widetilde{V}
=\int_{B^{4n}(\mathbf{0},\epsilon)}\left(u\circ\varphi_{\mathbf{a}}\cdot(\triangle \widehat{v})\circ\varphi_{\mathbf{a}}-\widehat{v}\circ\varphi_{\mathbf{a}}\cdot(\triangle u)\circ\varphi_{\mathbf{a}}\right) d\widetilde{V}\\
=&\int_{B^{4n}(\mathbf{a},\epsilon)}\left(u\circ\varphi_{\mathbf{a}}\cdot\triangle (\widehat{v}\circ\varphi_{\mathbf{a}})-\widehat{v}\circ\varphi_{\mathbf{a}}\cdot\triangle (u\circ\varphi_{\mathbf{a}})\right) d\widetilde{V}\\
=&2\int_{S^{4n-1}}\left.\left(\widehat{v}\circ\varphi_{\mathbf{a}}\cdot R^{1'}(u\circ\varphi_{\mathbf{a}})-u\circ\varphi_{\mathbf{a}}\cdot R^{0'}(\widehat{v}\circ\varphi_{\mathbf{a}})\right)\right|_{\epsilon\zeta}\frac{\epsilon^{4n-2}d\sigma(\zeta)}{(1-\epsilon^2)^{2n}}\\
=&2\int_{S^{4n-1}}\left.\left(G\cdot R^{1'}(u\circ\varphi_{\mathbf{a}})-u\circ\varphi_{\mathbf{a}}\cdot R^{0'}G\right)\right|_{\epsilon\zeta}\frac{\epsilon^{4n-2}d\sigma(\zeta)}{(1-\epsilon^2)^{2n}},
\end{aligned}\end{equation}
by taking the transformation $\mathbf{q}\rightarrow \varphi_{\mathbf{a}}(\mathbf{q})$, and using the invariance of $\triangle$ and the invariance of the measure $d\widetilde{V}$ under $\varphi_{\mathbf{a}}$. Here in the last identity, $\widehat{v}\circ\varphi_{\mathbf{a}}=G_{\mathbf{a}}\circ\varphi_{\mathbf{a}}=G$ on $S^{4n-1}(\mathbf{0},\epsilon)$.
Now substituting \eqref{part ii} into the second term in the right hand side of \eqref{total} and applying Green formula in Proposition \ref{Green} to the first term, we get
\begin{equation*}\begin{aligned}
\int_{S^{4n-1}}\left.(uR^{0'}G-G R^{1'}u)\right|_{r\zeta}\frac{r^{4n-2}d\sigma(\zeta)}{(1-r^2)^{2n}}
=\int_{S^{4n-1}}\left.\left(u\circ\varphi_{\mathbf{a}} R^{0'}G -G R^{1'}(u\circ\varphi_{\mathbf{a}})\right)\right|_{\epsilon\zeta}\frac{\epsilon^{4n-2}d\sigma(\zeta)}{(1-\epsilon^2)^{2n}}.
\end{aligned}\end{equation*}
Note that
\begin{equation*}\begin{aligned}
\int_{S^{4n-1}}\left.G\circ R^{1'}(u\circ\varphi_{\mathbf{a}})\right|_{\epsilon\zeta}\frac{\epsilon^{4n-2}d\sigma(\zeta)}{(1-\epsilon^2)^{2n}}\rightarrow 0,
\end{aligned}\end{equation*}
as $\epsilon\rightarrow 0$, by $g(\epsilon^2)\approx \epsilon^{2-4n}$ and $R^{1'}(u\circ\varphi_{\mathbf{a}})\rightarrow 0$, since coefficients of $R^{1'}$ tends to zero. Noting that
\begin{equation*}\begin{aligned}
R^{0'}G(\epsilon\zeta)=&g'(\epsilon^2)\epsilon^2=-\frac{(1-\epsilon^2)^{2n}}{\epsilon^{4n-2}},
\end{aligned}\end{equation*}
by \eqref{z jab}, we get
\begin{equation*}\begin{aligned}
&\int_{S^{4n-1}}\left.u\circ\varphi_{\mathbf{a}}\cdot R^{0'}G\right|_{\epsilon\zeta} \frac{\epsilon^{4n-2}d\sigma(\zeta)}{(1-\epsilon^2)^{2n}}\rightarrow-u(\mathbf{a}),
\end{aligned}\end{equation*}
as $\varepsilon\rightarrow 0$. Consequently,
$$u(\mathbf{a})=-\int_{S^{4n-1}}\left.(uR^{0'}G-GR^{1'}u)\right|_{r\zeta}\frac{r^{4n-2}d\sigma(\zeta)}{(1-r^2)^{2n}}.$$

  On the other hand, $G(\mathbf{q})\approx (1-|\mathbf{q}|^2)^{2n+1}$ as $\mathbf{q}\rightarrow 1$ by definition \eqref{Ga}, we have
\begin{equation*}\begin{aligned}
\int_{S^{4n-1}}\left.GR^{1'}u\right|_{r\zeta}\frac{r^{4n-2}d\sigma(\zeta)}{(1-r^2)^{2n}}\rightarrow0,
\end{aligned}\end{equation*}
as $r\rightarrow 1$. Thus
\begin{equation}\begin{aligned}\label{ua}
u(\mathbf{a})=-\lim_{r\rightarrow 1}\int_{S^{4n-1}}\left.uR^{0'}G\right|_{r\zeta}\frac{r^{4n-2}d\sigma(\zeta)}{(1-r^2)^{2n}}.
\end{aligned}\end{equation}
Since
\begin{equation}\begin{aligned}\label{psi eq}
1-|\varphi_{\mathbf{a}}(\mathbf{q})|^2=\frac{(1-|\mathbf{a}|^2)(1-|\mathbf{q}|^2)}{|1-\mathbf{a}^*\mathbf{q}|^2},
\end{aligned}\end{equation}
we have
\begin{align*}
R^{0'} G(\mathbf{q}) &= R^{0'} \left( g\left( |\varphi_{\mathbf{a}}(\mathbf{q})|^2 \right) \right) \\
&= g'(|\varphi_{\mathbf{a}}(\mathbf{q})|^2) \cdot R^{0'}\left( \frac{(1 - |\mathbf{a}|^2)(1 - |\mathbf{q}|^2)}{|1 - \mathbf{a}^* \mathbf{q}|^2} \right),\\
&= \left( \frac{(1 - |\mathbf{a}|^2)(1 - |\mathbf{q}|^2)}{|1 - \mathbf{a}^* \mathbf{q}|^2 - (1 - |\mathbf{a}|^2)(1 - |\mathbf{q}|^2)} \right)^{2n}   \left(\frac{-|\mathbf{q}|^2 (1 - |\mathbf{a}|^2)}{|1 - \mathbf{a}^* \mathbf{q}|^2}+O(1-|\mathbf{q}|^2)\right).
\end{align*}
Substitute this identity to \eqref{ua} to get
\begin{equation}\begin{aligned}\label{poisson ua}
u(\mathbf{a})
&= -\lim_{r \to 1} \int_{S^{4n-1}}u(r\zeta) \left( \frac{(1 - |\mathbf{a}|^2)(1 - r^2)}{|1 - \mathbf{a}^* r\zeta |^2 - (1 - |\mathbf{a}|^2)(1 - r^2)} \right)^{2n} \\
&\qquad\qquad\qquad\qquad\quad \cdot  \left( \frac{-r^2(1 - |\mathbf{a}|^2)}{|1 - \mathbf{a}^* r\zeta |^2}+O(1 - r^2) \right)\frac{r^{4n-2}}{(1 - r^2)^{2n}} d\sigma(\zeta) \\
&= \int u(\zeta) \cdot \left( \frac{1 - |\mathbf{a}|^2}{|1-\mathbf{a}^*\zeta|^2} \right)^{2n+1} d\sigma(\zeta).
\end{aligned}\end{equation}

For the uniqueness, let $\widetilde{u}$ be another solution to the Dirichlet problem \eqref{posson}. Apply the maximum principle in Theorem \ref{max} to $P[\varphi]-\widetilde{u}$, which is $\mathcal{QM}$-harmonic and equal to $0$ on the boundary. So we get $P[\varphi]-\widetilde{u}\equiv0$. The theorem is proved.
\end{proof}

\end{document}
